% Einheit 2 von 5 – Einsetzungsverfahren (bank/lineare-gleichungssysteme/e2.jsonl, zusammenbau v0.1)
%% TODO Zweigzeile Teil 1: was hier gelernt wird (Satz in Schülersprache aus dem Einheitstitel) – Einheit 2
\einheitenkopf[e2]{Einheit 2 von 5 · Einsetzungsverfahren}
\zweigzeile{ab Kl. 9 (am Gymnasium ab Kl. 8) · P10}
\setcounter{aufgabe}{15}

%% TODO Titel: Ich-kann-Satz fehlt in der Bank (Platzhalter: Kettenname) – Nr. 16
%% TODO Anweisung: ein Satz über der ersten Teilaufgabe fehlt in der Bank – Nr. 16
\begin{aufgabe}{Einsetzen}
\begin{teile}
\teil Welche Gleichung ist fast fertig? I: $3x + 7y = 17$, II: $x + 4y = 9$. Kreuze die Gleichung an, die du umstellst und in die andere einsetzt.\\ \kreuz{Gleichung I}\\ \kreuz{Gleichung II}
\end{teile}
%% TODO Darstellung für schwach fehlt in der Bank (lineare-gleichungssysteme-e2-k1-s1-v1; Abschnitt „Für schwache Schüler“)
\swz{$\text{I: $y = 2x + 1$, II: $3x + y = 16$ – Lösung? Setze I in II ein.}$}{}
%% TODO Darstellung für schwach fehlt in der Bank (lineare-gleichungssysteme-e2-k1-s1-v2; Abschnitt „Für schwache Schüler“)
\swz{$\text{I: $y = x + 3$, II: $2x + y = 15$ – Lösung? Setze I in II ein.}$}{}
%% TODO Darstellung für schwach fehlt in der Bank (lineare-gleichungssysteme-e2-k1-s1-v3; Abschnitt „Für schwache Schüler“)
\swz{$\text{I: $y = 3x - 2$, II: $x + y = 14$ – Lösung? Setze I in II ein.}$}{}
%% TODO Darstellung für schwach fehlt in der Bank (lineare-gleichungssysteme-e2-k1-s1-v4; Abschnitt „Für schwache Schüler“)
\swz{$\text{I: $y = 4x$, II: $x + 2y = 27$ – Lösung? Setze I in II ein.}$}{}
%% TODO Darstellung für schwach fehlt in der Bank (lineare-gleichungssysteme-e2-k1-s1-v5; Abschnitt „Für schwache Schüler“)
\swz{$\text{I: $y = 5x - 4$, II: $2x + y = 17$ – Lösung? Setze I in II ein.}$}{}
\swfrage{Was bleibt gleich, was ändert sich?}
%% TODO Darstellung für schwach fehlt in der Bank (lineare-gleichungssysteme-e2-k1-s2-v1; Abschnitt „Für schwache Schüler“)
\swz{$\text{I: $2x + y = 9$, II: $4x + y = 17$ – Lösung? Stelle I nach $y$ um und setze in II ein.}$}{}
%% TODO Darstellung für schwach fehlt in der Bank (lineare-gleichungssysteme-e2-k1-s3-v1; Abschnitt „Für schwache Schüler“)
\swz{$\text{I: $x + 2y = 7$, II: $x + 5y = 16$ – Lösung? Stelle I nach $x$ um und setze in II ein.}$}{}
%% TODO Darstellung für schwach fehlt in der Bank (lineare-gleichungssysteme-e2-k1-s4-v1; Abschnitt „Für schwache Schüler“)
\swz{$\text{I: $y = x + 2$, II: $3x + 4y = 29$ – Lösung? Setze I in II ein.}$}{}
%% TODO Darstellung für schwach fehlt in der Bank (lineare-gleichungssysteme-e2-k1-s5-v1; Abschnitt „Für schwache Schüler“)
\swz{$\text{I: $y = x + 4$, II: $5x - y = 16$ – Lösung? Setze I in II ein.}$}{}
%% TODO Darstellung für schwach fehlt in der Bank (lineare-gleichungssysteme-e2-k1-s6-v1; Abschnitt „Für schwache Schüler“)
\swz{$\text{I: $x + y = 3{,}50$, II: $4x + 5y = 15{,}10$ (Preise in Euro) – Lösung? Löse durch Einsetzen.}$}{}
%% TODO Darstellung für schwach fehlt in der Bank (lineare-gleichungssysteme-e2-k1-s7-v1; Abschnitt „Für schwache Schüler“)
\swz{$\text{I: $y = x + 5$, II: $2x + y = -1$ – Lösung? Löse durch Einsetzen.}$}{}
%% TODO Darstellung für schwach fehlt in der Bank (lineare-gleichungssysteme-e2-k1-s8-v1; Abschnitt „Für schwache Schüler“)
\swz{$\text{I: $y = 3x - 4$, II: $y = x + 2$ – Lösung? Setze die rechten Seiten gleich.}$}{}
%% TODO Darstellung für schwach fehlt in der Bank (lineare-gleichungssysteme-e2-k1-s9-v1; Abschnitt „Für schwache Schüler“)
\swz{$\text{I: $x + 3y = 14$, II: $2x - y = 7$ – Lösung mit Probe? Löse, mache die Probe in beiden Gleichungen und gib das Zahlenpaar an.}$}{}
\swz[3]{$\text{Fahrradladen: Ein Helm und ein Schloss kosten zusammen $54{,}00\,€$. Ein Verein kauft $7$ Helme und $4$ Schlösser für $316{,}20\,€$. Mit $x$ als Preis eines Helms und $y$ als Preis eines Schlosses in Euro gilt I: $x + y = 54$, II: $7x + 4y = 316{,}20$. Berechne die Preise und gib an, was ein Helm und was ein Schloss kostet. \hfill (P10 2022 OS)}$}{}
\end{aufgabe}

%% TODO Titel: Ich-kann-Satz fehlt in der Bank (Platzhalter: Kettenname) – Nr. 17
%% TODO Anweisung: ein Satz über der ersten Teilaufgabe fehlt in der Bank – Nr. 17
\begin{aufgabe}{Bruch als Lösung (Vorrat)}
%% TODO Darstellung für schwach fehlt in der Bank (lineare-gleichungssysteme-e2-k2-s1-v1; Abschnitt „Für schwache Schüler“)
\swz{$\text{I: $y = 2x$, II: $x + y = 5$ – Lösung? Löse durch Einsetzen, gib Brüche an.}$}{}
\end{aufgabe}

%% TODO Titel: Ich-kann-Satz fehlt in der Bank (Platzhalter: Kettenname) – Nr. 18
%% TODO Anweisung: ein Satz über der ersten Teilaufgabe fehlt in der Bank – Nr. 18
\begin{aufgabe}{Einsetzen – Fehler finden, Begründen}
\swz[3]{Paul löst I: $x + y = 6$, II: $2x + 3y = 14$: \rechnung{\text{I: } y &= 6 - x \\ \text{in I: } x + (6 - x) &= 6 \\ 6 &= 6} Er schreibt: „unendlich viele Lösungen“. Finde den Fehler.}{}
\swz[3]{Lisa löst I: $y = x + 3$, II: $4x + 2y = 30$: \rechnung{4x + 2(x + 3) &= 30 \\ 4x + 2x + 3 &= 30 \\ 6x &= 27} Finde den Fehler.}{}
\swz[3]{Jonas löst I: $y = 3x - 1$, II: $2x + y = 14$: \rechnung{2x + 3x - 1 &= 14 \\ 5x &= 15 \\ x &= 3} Er schreibt: „Lösung: 3“. Finde den Fehler.}{}
\swfrage{I: $y = 2x + 3$, II: $5x - y = 6$ – warum steht nach dem Einsetzen von I in II nur noch $x$ in der Gleichung? Begründe.}
\end{aufgabe}

\uebersichtskasten{Einsetzen: eine Gleichung nach einer Variablen umstellen und den Term in die andere Gleichung einsetzen – dann hat sie nur noch eine Variable. \\ I x + y = 9 \quad II 3x + 2y = 22 \quad I: y = 9 − x \quad in II: 3x + 2 · (9 − x) = 22 → 3x + 18 − 2x = 22 → x = 4 \quad y = 9 − 4 = 5 \\ Probe in beiden Gleichungen: 4 + 5 = 9 (wA) \quad 12 + 10 = 22 (wA) \quad Lösung (4 | 5). \\ Gleichsetzen: stehen beide Gleichungen in der Form y = …, die rechten Seiten gleichsetzen.}
